% Source: https://artofproblemsolving.com/community/c3h2303141_pigeonhole_principle
% Stable identity: aops-topic-2303141-q01
\textbf{(a)} Call a pair of cells in the same column a \emph{same-colour pair} if both cells have the same colour, and give it the \emph{type} (the two rows, the colour). There are $\dbinom52\cdot2=20$ possible types. If two different columns contain same-colour pairs of the same type, then those two rows and two columns form a monochromatic rectangle. So suppose, for contradiction, that all same-colour pairs in the grid have different types; then there are at most $20$ of them.

A column with $r$ red cells contains $\dbinom r2+\dbinom{5-r}{2}$ same-colour pairs. For $r=0,1,2,3,4,5$ this is $10,6,4,4,6,10$, so every column contains at least $4$, and the five columns contain at least $20$. Hence there are exactly $20$: every column has exactly $2$ or $3$ red cells, and every type occurs exactly once. In particular there are exactly $\dbinom52=10$ red same-colour pairs.

But a column with $3$ red cells has $3$ red pairs, and a column with $2$ red cells has $1$. If $t$ columns have $3$ red cells, the number of red pairs is $3t+(5-t)=5+2t$, which is odd and cannot equal $10$. This contradiction proves that a monochromatic rectangle always exists. $\blacksquare$

\textbf{(b)} Colour each of the $6$ columns with exactly two red cells, using each of the $\dbinom42=6$ possible pairs of rows exactly once:

\begin{asy}
usepackage("amsmath");
size(5cm);
defaultpen(fontsize(10pt));
int[][] rc={{1,1,1,0,0,0},{1,0,0,1,1,0},{0,1,0,1,0,1},{0,0,1,0,1,1}};
for(int i=0;i<4;++i) for(int j=0;j<6;++j){
  path sq=shift(j,-i)*unitsquare;
  if(rc[i][j]==1) filldraw(sq,pink,black); else filldraw(sq,lightblue,black);
  label(rc[i][j]==1?"R":"B",(j+0.5,-i+0.5),fontsize(9pt));
}
for(int j=0;j<6;++j) label("$"+string(j+1)+"$",(j+0.5,1),N,fontsize(8pt));
for(int i=0;i<4;++i) label("$"+string(i+1)+"$",(0,-i+0.5),W,fontsize(8pt));
\end{asy}

Suppose two columns and two rows had four red crossing cells. Each column has only two red cells, so in both columns the red cells would be exactly in those two rows, and the two columns would be coloured the same, which is not the case. The blue cells of a column are the other two rows, so the blue pairs of rows in the six columns are also all different, and the same argument rules out a blue rectangle. So this colouring has no monochromatic rectangle. $\blacksquare$
